\textbf{Phần 1. Trắc nghiệm lựa chọn}
\setcounterpageref{ex}{0}

\begin{ex} %[2H6N31]
Kim dài của đồng hồ quay 1 vòng (từ số 12 đến số 12) là được bao nhiêu phút?
\choice
{$30$ phút}
{$12$ phút}
{$24$ phút}
{\True $60$ phút (1 giờ)}
\loigiai{
Kim dài quay 1 vòng là 60 phút.\\
\textbf{Đáp án của câu hỏi là D}
}
\end{ex}

\begin{ex} %[2H6N31]
Khi kim ngắn chỉ số 6, kim dài chỉ số 12 thì đồng hồ chỉ mấy giờ?
\choice
{$12$ giờ}
{\True $6$ giờ đúng}
{$6$ giờ $30$ phút}
{$12$ giờ $30$ phút}
\loigiai{
Lúc đó là 6 giờ đúng.\\
\textbf{Đáp án của câu hỏi là B}
}
\end{ex}

\begin{ex} %[2H6H31]
Đồng hồ điện tử hiển thị $16:30$. Lúc đó là
\choice
{$6$ giờ $30$ phút chiều}
{$3$ giờ $30$ phút chiều}
{\True $4$ giờ $30$ phút chiều (hoặc 4 giờ rưỡi chiều)}
{$5$ giờ $30$ phút chiều}
\loigiai{
$16 - 12 = 4$ giờ 30 phút chiều.\\
\textbf{Đáp án của câu hỏi là C}
}
\end{ex}

\begin{ex} %[2H6H31]
Hôm nay là thứ Năm ngày 10. Hôm qua là
\choice
{\True Thứ Tư ngày 9}
{Thứ Ba ngày 9}
{Thứ Sáu ngày 11}
{Thứ Tư ngày 8}
\loigiai{
Hôm qua là thứ Tư ngày 9.\\
\textbf{Đáp án của câu hỏi là A}
}
\end{ex}

\begin{ex} %[2H6H31]
Hôm nay là thứ Bảy ngày 20. Ngày mai là
\choice
{Chủ nhật ngày 19}
{\True Chủ nhật ngày 21}
{Thứ Sáu ngày 19}
{Thứ Hai ngày 22}
\loigiai{
Ngày mai là Chủ nhật ngày 21.\\
\textbf{Đáp án của câu hỏi là B}
}
\end{ex}

\begin{ex} %[2H6H31]
Em đi ngủ lúc 21 giờ 30 phút. Lúc đó còn gọi là
\choice
{$8$ giờ $30$ phút tối}
{$10$ giờ $30$ phút tối}
{\True $9$ giờ $30$ phút tối}
{$9$ giờ tối}
\loigiai{
$21 - 12 = 9$ giờ 30 phút tối.\\
\textbf{Đáp án của câu hỏi là C}
}
\end{ex}

\begin{ex} %[2H6V31]
Bạn An vẽ tranh từ 8 giờ đến 8 giờ 30 phút. An đã vẽ tranh trong
\choice
{\True $30$ phút (nửa giờ)}
{$15$ phút}
{$45$ phút}
{$1$ giờ}
\loigiai{
Thời gian vẽ: 30 phút.\\
\textbf{Đáp án của câu hỏi là A}
}
\end{ex}

\begin{ex} %[2H6V31]
Bố đi làm lúc 7 giờ sáng và về nhà lúc 17 giờ. Bố đã ở nơi làm việc trong số giờ là
\choice
{$8$ giờ}
{$12$ giờ}
{$9$ giờ}
{\True $10$ giờ}
\loigiai{
$17 - 7 = 10$ giờ.\\
\textbf{Đáp án của câu hỏi là D}
}
\end{ex}

\begin{ex} %[2H6V31]
Tháng 10 có 31 ngày. Ngày cuối cùng của tháng 10 là ngày
\choice
{$30$}
{\True $31$}
{$28$}
{$29$}
\loigiai{
Ngày 31 tháng 10.\\
\textbf{Đáp án của câu hỏi là B}
}
\end{ex}

\begin{ex} %[2H6V31]
Một trận bóng đá bắt đầu lúc 15 giờ và kết thúc sau 90 phút (1 giờ 30 phút). Trận bóng kết thúc lúc
\choice
{\True $16$ giờ $30$ phút}
{$16$ giờ}
{$17$ giờ}
{$17$ giờ $30$ phút}
\loigiai{
$15$ giờ $+ 1$ giờ $30$ phút = $16$ giờ $30$ phút.\\
\textbf{Đáp án của câu hỏi là A}
}
\end{ex}

\begin{ex} %[2H6V31]
Gia đình Lan đi du lịch từ ngày 12 đến hết ngày 15 của tháng. Chuyến du lịch kéo dài bao nhiêu ngày?
\choice
{$3$ ngày}
{$5$ ngày}
{$6$ ngày}
{\True $4$ ngày (ngày 12, 13, 14, 15)}
\loigiai{
$15 - 12 + 1 = 4$ ngày.\\
\textbf{Đáp án của câu hỏi là D}
}
\end{ex}

\begin{ex} %[2H6C31]
Một đồng hồ quả lắc đánh chuông theo số giờ (1 giờ điểm 1 tiếng, 2 giờ điểm 2 tiếng...). Từ lúc 1 giờ đến 4 giờ chiều, đồng hồ điểm tất cả bao nhiêu tiếng chuông?
\choice
{$8$ tiếng}
{$9$ tiếng}
{\True $10$ tiếng ($1 + 2 + 3 + 4 = 10$)}
{$12$ tiếng}
\loigiai{
$1 + 2 + 3 + 4 = 10$ tiếng chuông.\\
\textbf{Đáp án của câu hỏi là C}
}
\end{ex}

\textbf{Phần 2. Trắc nghiệm đúng sai}
\setcounterpageref{ex}{0}

\begin{ex} %[2H6H31]
Xét tính đúng sai khi xem đồng hồ:
\choiceTF
{\True Kim ngắn chỉ giờ, kim dài chỉ phút.}
{\True Khi kim dài chỉ số 6 là 30 phút.}
{Khi kim dài chỉ số 12 là 12 phút.}
{\True Khi kim dài chỉ số 12 là 00 phút (giờ đúng).}
\loigiai{
\begin{itemize}
    \item Kim ngắn chỉ giờ, dài chỉ phút: Đúng.
    \item Chỉ số 6 là 30 phút: Đúng.
    \item Chỉ số 12 là 0 phút (không phải 12 phút). Ý 3 Sai, Ý 4 Đúng.
\end{itemize}
\begin{itemize}
    \item a) \textbf{Đúng}.
    \item b) \textbf{Đúng}.
    \item c) \textbf{Sai}.
    \item d) \textbf{Đúng}.
\end{itemize}
}
\end{ex}

\begin{ex} %[2H6V31]
Hôm nay là thứ Tư ngày 14. Xét tính đúng sai:
\choiceTF
{\True Hôm qua là thứ Ba ngày 13.}
{\True Ngày mai là thứ Năm ngày 15.}
{Thứ Tư tuần trước là ngày 8.}
{\True Thứ Tư tuần trước là ngày 7 ($14 - 7 = 7$).}
\loigiai{
\begin{itemize}
    \item Hôm qua: Thứ Ba ngày 13: Đúng.
    \item Ngày mai: Thứ Năm ngày 15: Đúng.
    \item $14 - 7 = 7 \neq 8$. Ý 3 Sai, Ý 4 Đúng.
\end{itemize}
\begin{itemize}
    \item a) \textbf{Đúng}.
    \item b) \textbf{Đúng}.
    \item c) \textbf{Sai}.
    \item d) \textbf{Đúng}.
\end{itemize}
}
\end{ex}

\textbf{Phần 3. Trả lời ngắn}
\setcounterpageref{ex}{0}

\begin{ex} %[2H6H31]
Khi kim ngắn chỉ số 7, kim dài chỉ số 12 thì đồng hồ chỉ mấy giờ đúng?
\shortans{7}
\loigiai{
\textbf{Đáp số:} $7$\\
7 giờ.
}
\end{ex}

\begin{ex} %[2H6H31]
Khi kim dài chỉ số 6 thì tương ứng với bao nhiêu phút?
\shortans{30}
\loigiai{
\textbf{Đáp số:} $30$\\
30 phút.
}
\end{ex}

\begin{ex} %[2H6V31]
Lúc 14 giờ 30 phút còn gọi là mấy giờ rưỡi chiều?
\shortans{2}
\loigiai{
\textbf{Đáp số:} $2$\\
$14 - 12 = 2$ giờ rưỡi chiều.
}
\end{ex}

\begin{ex} %[2H6V31]
Nếu thứ Sáu là ngày 8 thì thứ Sáu tuần sau là ngày bao nhiêu?
\shortans{15}
\loigiai{
\textbf{Đáp số:} $15$\\
$8 + 7 = 15$.
}
\end{ex}

\begin{ex} %[2H6V31]
Một tiết mục biểu diễn văn nghệ kéo dài từ 8 giờ đến 8 giờ 30 phút. Hỏi tiết mục kéo dài bao nhiêu phút?
\shortans{30}
\loigiai{
\textbf{Đáp số:} $30$\\
30 phút.
}
\end{ex}

\begin{ex} %[2H6C31]
Một tuần lễ em đi học 5 ngày, nghỉ 2 ngày cuối tuần. Hỏi trong 4 tuần lễ em được nghỉ bao nhiêu ngày cuối tuần?
\shortans{8}
\loigiai{
\textbf{Đáp số:} $8$\\
$2 \times 4 = 8$ ngày.
}
\end{ex}
